% Part: first-order-logic
% Chapter: introduction
% Section: first-order-logic

\documentclass[../../../include/open-logic-section]{subfiles}

\begin{document}

\olfileid{fol}{int}{fol}

\olsection{প্রথম-ক্রমের যুক্তিবিদ্যা}

বিধিবদ্ধ যুক্তিবিদ্যার প্রথম পরিচয় থেকেই সম্ভবত প্রথম-ক্রমের যুক্তিবিদ্যা
তোমার পরিচিত।\footnote{বস্তুত, আমরা মোটামুটি ধরেই নিচ্ছি যে তুমি এর সঙ্গে
পরিচিত!{} পরিচিত না হলে \emph{forall x} \citep{Magnus2021}-এর মতো আরও
প্রাথমিক কোনো পাঠ্যবই দেখে নিতে পারো।} একে তুমি “পরিমাণসূচকীয়
যুক্তিবিদ্যা” বা “বিধেয় যুক্তিবিদ্যা” নামেও চিনে থাকতে পারো। প্রথমত,
প্রথম-ক্রমের যুক্তিবিদ্যা একটি বিধিবদ্ধ ভাষা। অর্থাৎ, এর একটি নির্দিষ্ট
শব্দভাণ্ডার আছে এবং এর অভিব্যক্তিগুলি সেই শব্দভাণ্ডার দিয়ে গঠিত
প্রতীকক্রম। তবে প্রতিটি প্রতীকক্রম অনুমোদিত নয়। অনুমোদিত অভিব্যক্তিও নানা
ধরনের: পদ, !!{formula}s ও !!{sentence}s। প্রথম-ক্রমের যুক্তিবিদ্যার
!!{sentence}s নিয়েই আমাদের প্রধান আগ্রহ: এগুলি ইংরেজি বাক্যের একটি
বিধিবদ্ধ প্রতিরূপ দেয় এবং এগুলি সম্বন্ধে যুক্তিবিদদের স্বাভাবিক আগ্রহের
প্রশ্নগুলি করা যায়। যেমন:
\begin{itemize}
    \item $!B$ কি~$!A$ থেকে যৌক্তিকভাবে অনুসৃত হয়?
    \item $!A$ কি যৌক্তিকভাবে সত্য, যৌক্তিকভাবে মিথ্যা, না কি
আপতিক?
    \item $!A$ ও $!B$ কি সমতুল্য?
\end{itemize}

এই প্রশ্নগুলি মূলত প্রথম-ক্রমের যুক্তিবিদ্যার !!{sentence}s-এর “অর্থ”
সম্বন্ধে। যেমন, $!B$ কি~$!A$ থেকে যৌক্তিকভাবে অনুসৃত হয়—এই প্রশ্নটিকে
একজন দার্শনিক এভাবে বিশ্লেষণ করবেন: এমন কোনো ক্ষেত্র আছে কি যেখানে $!A$
সত্য কিন্তু~$!B$ মিথ্যা ($!B$, $!A$ থেকে অনুসৃত হয় না), না কি $!A$-কে
সত্য করে এমন প্রতিটি ক্ষেত্রই~$!B$-কেও সত্য করে ($!B$, $!A$ থেকে
অনুসৃত হয়)? কিন্তু “ক্ষেত্র” বলতে কী বোঝায়, তা এখনও বলা হয়নি—সেটিই
\emph{অর্থতত্ত্বের} কাজ। প্রথম-ক্রমের যুক্তিবিদ্যার অর্থতত্ত্ব দার্শনিকের
“ক্ষেত্র” সম্বন্ধীয় স্বজ্ঞাত ধারণার একটি গাণিতিকভাবে নির্ভুল মডেল দেয়;
আরও একটি গুরুত্বপূর্ণ বিষয়েরও মডেল দেয়—কোনো ক্ষেত্রে কোনো
!!a{sentence}~$!A$-র \emph{সত্য হওয়া} বলতে কী বোঝায়। আমরা যে
গাণিতিকভাবে নির্ভুল মডেলটি নির্মাণ করব, তাকে !!a{structure} বলি। “এর
মধ্যে সত্য” কথাটিকে নির্ভুল করে যে সম্পর্ক, তাকে \emph{পরিতৃপ্তি}-সম্পর্ক
বলা হয়। তাই !!{sentence}s~$!A$ ও !!{structure}s~$\Struct{M}$-এর জন্য
আমরা “$!A$,~$\Struct{M}$-এ পরিতৃপ্ত” (সংকেতে: $\Sat{M}{!A}$) সংজ্ঞায়িত
করব। তা হয়ে গেলে “থেকে অনুসৃত হয়” বা “যৌক্তিকভাবে সত্য”—এর মতো অন্য
অর্থগত পদগুলিরও নির্ভুল সংজ্ঞা দিতে পারব। তখন, যেমন,
$\lforall[x][(!A(x) \lif !B(x))], \lexists[x][!A(x)] \Entails
\lexists[x][!B(x)]$ কি না, আবারও গাণিতিক নির্ভুলতায় নির্ধারণ করা সম্ভব
হবে। উত্তরটি অবশ্যই “হ্যাঁ”। ইংরেজি বাক্যকে প্রথম-ক্রমের যুক্তিবিদ্যায়
সংকেতায়িত করার প্রশিক্ষণ আগে পেয়ে থাকলে একে তুমি, যেমন, “সব পিঁপড়ে
কীট, পিঁপড়ে আছে, অতএব কীট আছে”—এই যুক্তির সংকেতায়ন বলে চিনতে পারবে।
যুক্তিটি স্পষ্টতই বৈধ; তাই আমাদের বিধিবদ্ধ ভাষায় “থেকে অনুসৃত হয়”-এর
গাণিতিক মডেলটিরও একই উত্তর দেওয়া উচিত।% BN-SRC-083 TARGET-CORRECTION-BEGIN
\par\smallskip\noindent\textnormal{\small\textbf{উৎস-সংশোধন (BN-SRC-083):} হিমায়িত ইংরেজি উৎসের প্রদর্শিত অনুসিদ্ধান্তে সার্বিক সূত্রটির বন্ধনী শেষ না করে পরের পূর্বধারণা ও অনুসিদ্ধান্ত-সংকেতকে তার পরিসরের মধ্যে ঢুকিয়ে দেওয়া হয়েছে এবং একেবারে শেষে অতিরিক্ত বন্ধনী আছে। বাংলা পাঠে সার্বিক সূত্রটির পরেই বন্ধনী শেষ করে দুই পূর্বধারণা ও সিদ্ধান্তকে পৃথক রাখা হয়েছে।} % BN-SRC-083 TARGET-CORRECTION-END

বিধিবদ্ধ যুক্তিবিদ্যার প্রথম পরিচয় থেকে সম্ভবত তোমার মনে আছে যে
\emph{!!{derivation}s}-ও আছে। প্রথম কোনো বিধিবদ্ধ যুক্তিবিদ্যার পাঠ নিয়ে
থাকলে শিক্ষক নিশ্চয়ই তোমাকে এমন !!{derivation}s খোঁজার অনুশীলন
করিয়েছেন; এমনকি উপরের অনুসিদ্ধান্তটি যে সিদ্ধ হয়, তা দেখায় এমন
!!a{derivation}-ও হয়তো। !!{derivation}s দেওয়ার নানা উপায় আছে: তুমি
হয়তো “স্বাভাবিক নিষ্পাদন” বা “সত্য-বৃক্ষ” নামে কোনো পদ্ধতি ব্যবহার করেছ,
তবে আরও অনেক পদ্ধতি আছে। !!{derivation} পদ্ধতির উদ্দেশ্য হলো এমন উপকরণ
দেওয়া, যার সাহায্যে যুক্তিবিদের উপরের প্রশ্নগুলির উত্তর দেওয়া যায়।
যেমন, এমন একটি স্বাভাবিক নিষ্পাদন !!{derivation}, যেখানে
$\lforall[x][(!A(x) \lif !B(x))]$ ও $\lexists[x][!A(x)]$ পূর্বধারণা এবং
$\lexists[x][!B(x)]$ সিদ্ধান্ত (শেষ পংক্তি), সেটি \emph{যাচাই করে} যে
$\lexists[x][!B(x)]$, $\lforall[x][(!A(x) \lif !B(x))]$ ও
$\lexists[x][!A(x)]$ থেকে যৌক্তিকভাবে অনুসৃত হয়।

কিন্তু কেন? আপাতদৃষ্টিতে !!{derivation} পদ্ধতির সঙ্গে অর্থতত্ত্বের কোনো
সম্পর্কই নেই: কোনো বিধিবদ্ধ !!{derivation} দেওয়া মানে কেবল কয়েকটি
বিধিনিয়ন্ত্রিত উপায়ে সংকেত সাজানো; সেখানে “ক্ষেত্র” বা “এর মধ্যে সত্য”
কিছুরই উল্লেখ নেই। !!{derivation} পদ্ধতি ও অর্থতত্ত্বের সংযোগটি একটি
অধিযৌক্তিক অনুসন্ধানের মাধ্যমে প্রতিষ্ঠা করতে হয়। যেমন, একটি গাণিতিক
প্রমাণ দরকার যে $\lforall[x][(!A(x) \lif !B(x))]$ ও
$\lexists[x][!A(x)]$ পূর্বধারণা থেকে $\lexists[x][!B(x)]$-এর একটি
বিধিবদ্ধ !!{derivation} সম্ভব, যদি এবং কেবল যদি
$\lforall[x][(!A(x) \lif !B(x))]$ ও $\lexists[x][!A(x)]$ একত্রে
$\lexists[x][!B(x)]$-কে অনুসিদ্ধ করে। কিন্তু তা করার আগে সংকেতবিন্যাস ও
অর্থতত্ত্বের সংজ্ঞাগুলি ঠিক করতে অনেক শ্রমসাধ্য কাজ সম্পন্ন করতে হবে।% BN-SRC-084 TARGET-CORRECTION-BEGIN
\par\smallskip\noindent\textnormal{\small\textbf{উৎস-সংশোধন (BN-SRC-084):} হিমায়িত ইংরেজি উৎসে স্বাভাবিক নিষ্পাদনের উদাহরণ ও তার অধিযৌক্তিক পুনর্বিবৃতিতে সার্বিক পূর্বধারণার শেষ বন্ধনী দুবার বাদ গেছে এবং অস্তিত্বমূলক সিদ্ধান্তের শেষে দুবার অতিরিক্ত বন্ধনী আছে। বাংলা পাঠে চারটি সূত্রের বন্ধনী তাদের নিজ নিজ পরিসর অনুসারে সুষম করা হয়েছে।} % BN-SRC-084 TARGET-CORRECTION-END

\end{document}
